\begingroup\fontsize{9.25}{11.4}\selectfont
\setlength{\tabcolsep}{3pt}\renewcommand{\arraystretch}{1.18}
\begin{longtable}{@{}>{\raggedright\arraybackslash}p{0.38\linewidth}rp{0.16\linewidth}p{0.26\linewidth}@{}}
\caption{Project composition and assignment jumps through 2016}\label{tab:pub_A08} \\
\toprule Project group & Records / CCPP & Share (\%) & Assignment jump (SE), pp \\ \midrule\endfirsthead
\multicolumn{4}{c}{\tablename\ \thetable{} -- continued} \\
\toprule Project group & Records / CCPP & Share (\%) & Assignment jump (SE), pp \\ \midrule\endhead
Productive livelihood & 96 / 96 & 54.5 & 60.94 (26.11) \\
Social and basic services & 31 / 31 & 17.6 & 34.73 (26.39) \\
Community and civic infrastructure & 36 / 36 & 20.5 & 17.30 (10.68) \\
Management and capacity support & 13 / 13 & 7.4 & 0.21 (0.62) \\
\bottomrule\end{longtable}
\parbox{0.97\linewidth}{\footnotesize \textit{Notes:} Composition covers CMAN project records through 2016 linked to the 389-community complete Census 2017 CCPP analysis universe. A community may receive projects in more than one group. Cutoff jumps are robust bias-corrected local-linear assignment discontinuities in receipt of at least one group-specific project, using triangular weights, \(h=0.0075\), \(b=0.0135\), mass-point adjustment, and district CR2 inference. Project type is a post-assignment implementation attribute; the table does not estimate causal project-type heterogeneity. Sources: RUV, CMAN, INEI 2007 Census, and INEI-assisted Census 2017.}
\par\smallskip{\footnotesize\textit{Review boundary:} Conditional local evidence; historical design search and selection remain limitations. RF denotes the assignment reduced form; CCPP denotes a community. In outcome tables, shares, binary outcomes and zero-to-one indices are in percentage points; logarithms remain log points and counts remain counts. A common window is not claimed optimal for every outcome. No artifact-level release approval or manuscript synchronization is implied.}\par\endgroup
